\documentclass[11pt,a4paper]{article}
\usepackage[margin=25mm]{geometry}
\usepackage{amsmath,amssymb,booktabs,array,longtable}
\usepackage[hidelinks]{hyperref}
\usepackage{microtype}

\title{Cutoff Functions Implemented in AccelNet}
\author{AccelNet development notes}
\date{\today}

\newcommand{\rc}{R_{\mathrm c}}
\newcommand{\fc}{f_{\mathrm c}}

\begin{document}
\maketitle

\section{Conventions}

For cutoff types 1 and 4--8, define the inner radius and normalized transition
coordinate by
\begin{equation}
  r_{\mathrm{in}}=\alpha\rc,
  \qquad
  x=\frac{r-r_{\mathrm{in}}}{\rc-r_{\mathrm{in}}},
  \qquad 0\leq\alpha<1.
\end{equation}
Their common piecewise form is
\begin{equation}
  \fc(r)=
  \begin{cases}
    1, & 0\leq r<r_{\mathrm{in}},\\
    F(x), & r_{\mathrm{in}}\leq r<\rc,\\
    0, & r\geq\rc,
  \end{cases}
  \qquad
  \frac{d\fc}{dr}=
  \begin{cases}
    0, & 0\leq r<r_{\mathrm{in}},\\
    \dfrac{F'(x)}{\rc-r_{\mathrm{in}}}, & r_{\mathrm{in}}\leq r<\rc,\\
    0, & r\geq\rc.
  \end{cases}
\end{equation}
The hard and tanh cutoffs ignore $\alpha$.  Type 9 instead interprets
$\alpha=h/\rc$ and requires $0<\alpha<1$.

\section{Implemented functions}

\subsection*{Type 0: hard cutoff}
\begin{equation}
  \fc(r)=\begin{cases}1,&r\leq\rc,\\0,&r>\rc,\end{cases}
  \qquad \frac{d\fc}{dr}=0\quad(r\neq\rc).
\end{equation}
This cutoff is discontinuous at $\rc$; the implementation returns zero for
its ordinary derivative at the boundary.

\subsection*{Type 1: cosine cutoff}
\begin{equation}
  F(x)=\frac{1+\cos(\pi x)}{2},
  \qquad F'(x)=-\frac{\pi}{2}\sin(\pi x).
\end{equation}

\subsection*{Types 2 and 3: tanh cutoffs}
Let $t=\tanh(1-r/\rc)$.  For $r<\rc$,
\begin{align}
  \text{type 2:}\quad &\fc(r)=t^3,
  &\frac{d\fc}{dr}=\frac{3t^2(t^2-1)}{\rc},\\
  \text{type 3:}\quad &\fc(r)=\frac{t^3}{\tanh^3(1)},
  &\frac{d\fc}{dr}=\frac{3t^2(t^2-1)}{\rc\tanh^3(1)}.
\end{align}
For $r\geq\rc$, both the value and derivative are zero.  Type 3 is normalized
so that $\fc(0)=1$.

\subsection*{Type 4: exponential cutoff}
\begin{equation}
  F(x)=\exp\!\left(1+\frac{1}{x^2-1}\right),
  \qquad
  F'(x)=-\frac{2x}{(x^2-1)^2}
  \exp\!\left(1+\frac{1}{x^2-1}\right).
\end{equation}

\subsection*{Types 5--8: polynomial cutoffs}
\small
\begin{longtable}{c>{$}l<{$}>{$}l<{$}}
\toprule
Type & F(x) & F'(x)\\
\midrule
5 & 1-3x^2+2x^3 & -6x+6x^2\\
6 & 1-10x^3+15x^4-6x^5 & -30x^2+60x^3-30x^4\\
7 & 1-35x^4+84x^5-70x^6+20x^7
  & -140x^3+420x^4-420x^5+140x^6\\
8 & 1-126x^5+420x^6-540x^7+315x^8-70x^9
  & -630x^4+2520x^5-3780x^6+2520x^7-630x^8\\
\bottomrule
\end{longtable}
\normalsize
Type $4+m$ has its first $m$ derivatives equal to zero at the transition
endpoints, for $m=1,\ldots,4$.

\subsection*{Type 9: fractional cutoff}
The fractional cutoff introduced in Eq.~(A5) of Mori \emph{et al.} is
\begin{equation}
  X=\frac{r-\rc}{h}=\frac{r-\rc}{\alpha\rc},
  \qquad
  \fc(r)=
  \begin{cases}
    \dfrac{X^2}{1+X^2},&r<\rc,\\
    0,&r\geq\rc,
  \end{cases}
\end{equation}
with analytical derivative
\begin{equation}
  \frac{d\fc}{dr}=
  \begin{cases}
    \dfrac{2X}{\alpha\rc(1+X^2)^2},&r<\rc,\\
    0,&r\geq\rc.
  \end{cases}
\end{equation}
The paper uses $\rc=0.45\,\mathrm{nm}$ and $h=0.125\,\mathrm{nm}$ for its
three-body contribution.  This corresponds to
$\alpha=h/\rc=5/18\simeq0.2777778$ in AccelNet.

\section{Use by descriptor family}

The same selected cutoff is applied consistently to values and analytical
derivatives:
\begin{itemize}
  \item Behler radial functions contain one factor $\fc(r_{ij})$; angular
        functions contain the appropriate products over their radial legs.
  \item Chebyshev radial terms contain $\fc(r_{ij})$, while angular pair terms
        contain $\fc(r_{ij})\fc(r_{ik})$.  Direct and Cartesian-moment angular
        evaluation use the same factors.
  \item LJ descriptors use $r^{-6}\fc(r)$ and $r^{-12}\fc(r)$, including the
        product rule in their Cartesian derivatives.
\end{itemize}

In setup files, use for example
\begin{verbatim}
BASIS type=Chebyshev cutoff_type=9 cutoff_alpha=0.2777777777777778
BASIS type=LJ        cutoff_type=9 cutoff_alpha=0.2777777777777778
FUNCTIONS type=Behler2011 cutoff_type=9 cutoff_alpha=0.2777777777777778
\end{verbatim}
Legacy defaults are type 1 for Chebyshev and Behler, and type 0 for LJ.
Extended Aenet network files store the cutoff type and alpha in descriptor
parameter rows 5 and 6.

\section{Reference}

H. Mori \emph{et al.}, ``Dynamic interaction between dislocations and
obstacles in bcc iron based on atomic potentials derived using neural
networks,'' \emph{Physical Review Materials} \textbf{7}, 063605 (2023),
Appendix A, Eqs.~(A4)--(A6), DOI:
\href{https://doi.org/10.1103/PhysRevMaterials.7.063605}
{10.1103/PhysRevMaterials.7.063605}.

\end{document}
